\documentclass{yijiem}
\usepackage[colorlinks=true, linkcolor=blue, urlcolor=blue]{hyperref}
\begin{document}

\begin{spacing}{1.1}
\name{Yijie Ma}
\end{spacing}
\contact{605 E Evelyn Ave, Mountain View, CA 94041}{yijiem.main@gmail.com}{\href{https://github.com/yijiem}{github}}{650-695-9643}

\section{Skills}
    C/C++, Go, Python, Distributed Systems, AI Infrastructure (GenAI, MCP, A2A, GPU/RDMA), gRPC, API Gateway/L7 Proxy, Performance Optimization, OpenTelemetry, Storage \& Filesystems (RocksDB), \mbox{Kubernetes}, Networking

\section{Work Experience}
\datedsubsection{Google}{April 2017 - Present}
    \textit{\small Software Engineer}
    \begin{description}\setlength{\topsep}{2pt}\setlength{\partopsep}{0pt}\setlength{\itemsep}{1.5pt}\setlength{\parskip}{0pt}
    \projectitem{Google Cloud API Gateway \& Service Infrastructure}{\workitemsfive
        {Led technical direction and C++ distributed tracing architecture for Gemini Agent Platform across HTTP/gRPC, defining OTel SemConv data models for Cloud, GenAI, \href{https://docs.cloud.google.com/mcp/monitor-mcp-tool-use-with-cloud-trace}{Model Context Protocol (MCP)}, and Agent-to-Agent (A2A).}
        {Drove cross-org alignment across Auth, Observability, and API Gateway teams to resolve ambiguous agentic requirements; scoped milestones and mentored junior engineers on multi-quarter observability projects.}
        {Built end-to-end W3C Trace Context propagation, trace span generation, trace context initiation and sampling, and asynchronous OTLP trace/log export pipelines for a C++ L7 proxy handling millions of QPS.}
        {Scoped and led multi-quarter delivery of an isolated fast API serving release track across control-plane and serving teams, decoupling rapid AI and Cloud feature rollouts from the slower stable release track.}
        {Pioneered LLM multi-agent workflows using reusable AI Skills to automate release risk evaluation, qualification reporting, and staging promotions, adopted across the team.}}
    \projectitem{gRPC Open Source Remote Procedure Call Framework}{\workitemsfour
        {Core developer and \href{https://github.com/grpc/grpc/blob/master/MAINTAINERS.md}{Maintainer} of \href{https://github.com/grpc/grpc}{gRPC} (40K+ GitHub stars) with 150+ authored PRs, driving open-source RFCs and PR reviews across C++ core, networking, and observability.}
        {Co-authored \href{https://github.com/grpc/proposal/blob/master/A79-non-per-call-metrics-architecture.md}{gRFC A79} to standardize non-per-call telemetry; architected the C++ templatized \href{https://github.com/grpc/grpc/pull/36449}{Metrics framework} and \href{https://github.com/grpc/grpc/pull/36070}{OpenTelemetry plugin}, \href{https://github.com/grpc/grpc/pull/37311}{cutting Counter overhead by 10\%}.}
        {Designed and built the cross-platform EventEngine asynchronous \href{https://github.com/grpc/grpc/pull/32701}{DNS resolver} across Linux/POSIX and Windows.}
        {Designed POSIX fork syscall handling and resolved multi-threaded pre-fork deadlocks across gRPC C++ and Python.}}
    \projectitem{Google Kubernetes Engine (GKE) Networking}{\workitemstwo
        {Engineered GKE control-plane automation in Go to auto-provision multi-NIC \href{https://cloud.google.com/ai-hypercomputer/docs/create/gke-ai-hypercompute-custom\#create-cluster}{RDMA network topologies} (VPCs, subnets, IP ranges) for AI-optimized GPU clusters.}
        {Designed concurrent request handling with distributed transactions and validated multi-node GPU communication using NVIDIA's NCCL test suite.}}
    \projectitem{Google Drive Desktop Application}{\workitemsfour
        {Designed and built the \href{https://workspaceupdates.googleblog.com/2025/01/release-notes-01-10-2025.html}{rsync-style differential upload} and Gzip stream compression encoders with compressibility heuristics to save upload bandwidth for millions of users.}
        {Architected the \href{https://blog.google/products-and-platforms/products/workspace/drive-for-desktop/}{USB device backup subsystem} backing up millions of files daily, increasing backup speed by 10x (140s to 18s for 1GiB) by optimizing I/O read buffers.}
        {Added Crashpad crash reporting on Linux, built a post-mortem crash debugging tool, and resolved a new category of Windows crashes (1000+ daily) in threads running on heap-allocated stacks.}
        {Designed the \href{https://knowledge.workspace.google.com/admin/drive/advanced-drive-for-desktop-configuration}{cross-platform Preferences system} supporting user-level settings and enterprise admin policy overrides across macOS and Windows.}}
    \projectitem{Stadia Cloud Gaming Service}{\workitemstwo
        {Increased game click-to-play launch speed by 75\% (1.58s to 396ms) via CPU profiling and container setup optimization.}
        {Designed a Linux Cgroup library for isolating and configuring container resources on cloud gaming VMs.}}
    \end{description}
\datedsubsection{NetApp}{August 2015 - March 2017}
    \textit{\small Member of Technical Staff}\par\vspace{1pt}
    Worked on the \href{https://www.netapp.com/media/19930-ds-3651.pdf}{AltaVault} Cloud-Integrated Storage device. Designed and built a filesystem software to backup user's files to multi-cloud vendors. Designed and implemented full and incremental snapshot backup and restore functions using RocksDB.

\section{Education}
\datedsubsection{Carnegie Mellon University}{Sep 2013 - May 2015}
    \begin{spacing}{0.7}
    \textit{\small Master of Science in Electrical and Computer Engineering}\newline
    \end{spacing}
    \textbf{Coursework:} Introduction to Computer System, Embedded Realtime System, Advanced Storage System, Internet Service, Cloud Computing, Distributed System
\datedsubsection{Donghua University}{Sep 2009 - May 2013}
    \begin{spacing}{0.7}
    \textit{\small Bachelor of Engineering in Electrical Engineering and Automation}\newline
    \end{spacing}
    \textbf{Awards:} 2012 IEEE-standard Micromouse competition national award issued by China Computer Federation (CCF)
\end{document}
